\documentclass[11pt]{article}

\usepackage[margin=1in]{geometry}
\usepackage{enumitem}
\usepackage{amsmath, amssymb}
\usepackage{hyperref}

\setlist[enumerate]{itemsep=0.8em, topsep=0.8em}
\setlist[itemize]{itemsep=0.3em, topsep=0.3em}

\begin{document}

\begin{center}
  {\Large \textbf{PHYS 455 Examination}}\\[0.25em]
  {\large \textbf{Quantum Mechanics and Wave Functions}}\\[0.5em]
\end{center}

\noindent\textbf{Instructions:}
\begin{itemize}
  \item Answer all questions.
  \item For Questions 1--4, choose the best option.
  \item For Questions 5--8, write brief but precise answers.
\end{itemize}

\vspace{0.5em}

\begin{enumerate}[label=\textbf{\arabic*.}]

  \item In the context of the uncertainty principle, what is the minimum value of the product 
  $\Delta x \cdot \Delta p$ for a quantum particle?
  \begin{enumerate}[label=(\Alph*)]
    \item $\hbar$
    \item $\hbar/2$
    \item $h$
    \item $h/4\pi$
  \end{enumerate}

  \item For a particle in a one-dimensional infinite square well of width $L$, which statement 
  correctly describes the energy eigenvalues?
  \begin{enumerate}[label=(\Alph*)]
    \item They are evenly spaced with separation $\hbar^2/2mL^2$
    \item They scale as $n^2$ where $n$ is the quantum number
    \item They are independent of the well width $L$
    \item They form a continuous spectrum
  \end{enumerate}

  \item What physical interpretation does the Born rule provide for the wave function $\psi(x,t)$ 
  of a quantum particle?
  \begin{enumerate}[label=(\Alph*)]
    \item $|\psi(x,t)|$ represents the probability density
    \item $\psi(x,t)$ represents the particle's definite position
    \item $|\psi(x,t)|^2$ represents the probability density for finding the particle at position $x$
    \item $\text{Re}(\psi(x,t))$ represents the observable position
  \end{enumerate}

  \item Which property distinguishes a Hermitian operator in quantum mechanics and makes it suitable 
  for representing physical observables?
  \begin{enumerate}[label=(\Alph*)]
    \item It always commutes with the Hamiltonian
    \item It has real eigenvalues and orthogonal eigenfunctions
    \item It preserves the norm of wave functions under time evolution
    \item It generates unitary transformations exclusively
  \end{enumerate}

  \item Explain the physical significance of commutators in quantum mechanics. If two observables 
  have a non-zero commutator, what does this imply about simultaneous measurements of these observables?

  \item Describe the quantum harmonic oscillator and explain why it serves as a fundamental model 
  in quantum mechanics. What are the key features of its energy spectrum and wave functions?

  \item State the time-dependent Schr\"odinger equation and explain its role as the fundamental 
  equation of motion in quantum mechanics. How does it differ from the time-independent version, 
  and when is each appropriate?

  \item Explain the concept of quantum tunneling. How does it arise from the wave nature of particles, 
  and what are its implications for particles encountering potential barriers that would be classically 
  insurmountable?

\end{enumerate}

\end{document}
